\documentclass[conference]{IEEEtran}
\usepackage[utf8]{inputenc}
\usepackage[T1]{fontenc}
\usepackage{cite}
\usepackage{amsmath,amssymb}
\usepackage{graphicx}
\usepackage{booktabs}
\usepackage{url}

\title{A Preregistered Audit of Dataset Contamination Controls in LLM-for-NetOps Benchmarks}

\author{
\IEEEauthorblockN{Autonomous AI Researcher}
\IEEEauthorblockA{CNSM 2026 Agentic AI Researcher Track}
}

\begin{document}

\maketitle

\begin{abstract}
We preregistered and executed a paired, confirmatory audit to test whether conservative deterministic contamination detectors (exact‑match + fuzzy 5‑gram + provenance heuristics) flag items in a reproducible 150‑item NetOps sample and whether applying preregistered contamination controls changes validator‑based success rates. Crucially, the executed sample used provenance‑stable synthetic tasks (source\_identifier prefix "synthetic-netops:") and shared generation semantics (independent\_condition\_generation = false), recorded in the archived model and task manifests; these two execution facts materially limited the audit's sensitivity to generation‑level contamination. No preregistered high‑confidence contamination flags occurred and therefore no guarded exclusions or replacements were performed. All 150 baseline and matched guarded episodes passed the deterministic validator (baseline = 150/150, guarded = 150/150). The paired success‑rate difference = 0.0 percentage points (95\% CI [0.0, 0.0]); exact McNemar two‑sided p = 1.0. We reconcile the preregistration\ensuremath{\rightarrow}execution gap with artifact‑grounded diagnosis, explain why the run is degenerate for detecting public‑benchmark leakage, and provide concrete, artifact‑anchored recommendations for audit designs that would be sensitive to contamination in web‑sourced benchmarks.
\end{abstract}

\section{Introduction}
Motivation. Large language models (LLMs) are increasingly applied to NetOps tasks such as intent‑to‑configuration translation and automated configuration repair. Operational decisions based on benchmark results require that measured performance reflect model capability rather than dataset contamination or templated item leakage into model training. Meta‑analyses and recent surveys call for contamination‑aware evaluations and validator‑backed oracles to avoid inflated reliability claims.

Research question and contribution. We preregistered and executed an artifact‑anchored audit to answer: under conservative deterministic contamination detectors (exact normalized token equality; fuzzy 5‑gram shingle overlap with preregistered thresholds; provenance timestamp heuristics) and a guarded exclusion policy, how many high‑confidence flags arise in a reproducible 150‑item NetOps sample, and how much does applying those controls change a single hosted LLM's validator pass rate? Our contributions are: (1) a fully preregistered confirmatory experiment (study\_id = contamination-audit-netops-2026-v1) executed end‑to‑end and archived; (2) exact, artifact‑verified confirmatory statistics on Npairs = 150 with deterministic validator traces; (3) an explicit preregistration\ensuremath{\rightarrow}execution reconciliation that identifies why the run produced a degenerate null and prescribes artifact‑grounded design changes to increase audit sensitivity.

Scope and bounded claims. The audit intentionally targeted a methodological question about preregistered contamination controls, not an empirical prevalence estimate for public web‑sourced benchmarks. The executed confirmatory sample used provenance‑stable synthetic tasks (source\_identifier prefix "synthetic-netops:") produced by the deterministic task generator and employed shared\_initial\_candidate generation semantics (independent\_condition\_generation = false). Because both facts are recorded in the archived model and task manifests, they limit inference: the observed zero‑flag, zero‑difference outcome applies to this synthetic sampling and generation regime and does not imply absence of contamination in independent public benchmarks.

\section{Related Work}
LLM‑driven NetOps systems and verifier integration. Recent system and evaluation work has investigated LLMs for intent‑driven configuration synthesis and for proposing configuration repairs; these studies highlight practical value but also recurring failure modes that motivated tool‑assisted or agentic architectures. Representative intent\ensuremath{\rightarrow}configuration systems show that LLMs can produce syntactically valid device artifacts when guided by structured prompts or in‑context examples, but that semantic correctness (reachability, policy preservation) typically requires an external validator or iterative refine‑and‑repair loop [https://openalex.org/W4404602270, 10.1109/vtc2025-fall65116.2025.11310185]. Benchmarks and controlled experiments on LLM‑assisted repair report improved average repair quality when the pipeline integrates a deterministic checker (or formal verifier) with generation and repair steps; however, per‑case regressions remain common and depend strongly on the verifier's property coverage and on dataset realism [https://openalex.org/W7157506260, https://openalex.org/W7163636922].

Formal verification and deterministic oracles. The NetOps literature provides mature deterministic verification techniques (Header Space Analysis and derived tools, localized subspecification approaches, and NetKAT/related formalizations) that are practical oracles for reachability and many policy properties; these tools form the empirical basis for validator‑backed evaluation in NetOps generate‑validate pipelines [http://citeseerx.ist.psu.edu/viewdoc/summary?doi=10.1.1.300.8659, https://openalex.org/W2742515467, https://openalex.org/W1882012874]. Studies integrating such verifiers show that the choice of checks (which fields and temporal/workflow constraints are modeled) materially changes which generation errors are visible to the auditor---an observation we incorporate by using a deterministic validator (deterministic\_netops\_validator\_v5) that emits parse\ensuremath{\rightarrow}state traces and explicit violation codes for every episode.

Agentic validate--repair architectures and their limits. Multiple contemporaneous projects evaluate agentic or multi‑stage designs where an LLM proposes a candidate, a verifier checks it, and a repair generator attempts fixes when violations are found; these generate--validate--repair patterns generally improve mean repair rates versus single‑pass generation but do not eliminate regressions and depend on the repair‑policy and localization strategy [https://openalex.org/W7162218766, https://openalex.org/W7163636922]. Empirical studies also show scale sensitivity: effectiveness reported on curated or synthetic misconfiguration tasks often degrades when confronted with more heterogeneous, multi‑vendor, or production‑sourced corpora. These findings motivate two audit recommendations we make below: (1) use provenance‑rich sampling to exercise contamination detectors, and (2) persist richer validator/tracer outputs to analyze subtle regression modes.

Benchmark validity, contamination, and detector strategies. Recent surveys and meta‑evaluations emphasize two recurring evaluation validity risks: dataset contamination, where benchmark items (or templated prompts) are present in LLM training data and inflate measured accuracy, and construct invalidity, where tasks are overly templated or low complexity and fail to exercise operational difficulty [https://openalex.org/W4404534210, https://openalex.org/W7166900660]. The literature recommends deterministic fingerprinting (canonicalized exact matching), fuzzy n‑gram overlap, and provenance/timestamp tracing as complementary heuristic detectors; these techniques form the basis for many proposed audit workflows because they are computationally tractable and interpretable. However, prior work also highlights tradeoffs: exact fingerprints miss paraphrased memorization; n‑gram overlap thresholds require sensitivity analysis; and provenance heuristics depend on available source metadata and archival coverage. Our preregistration implemented a conservative combination (exact normalized equality OR fuzzy 5‑gram overlap >= 0.80 for high confidence, with lower thresholds for sensitivity analyses) precisely because the literature documents these tradeoffs and recommends multimodal flagging and post‑hoc adjudication [https://openalex.org/W4385681611, https://openalex.org/W7166900660].

How our study differs from and complements prior work. Unlike broader capability benchmarks that measure absolute LLM performance across many models, or agentic repair evaluations that emphasize mean repair rates, our study is a preregistered contamination‑control audit: it asks whether a conservative, preregistered detector ensemble would produce high‑confidence flags in a reproducible 150‑item sample and whether deterministic guarded exclusions would alter validator pass rates. Prior benchmark and system papers typically do not publish preregistered exclusion rules or the exact generation semantics used for paired comparisons; that omission can obscure whether observed generation discordance arises from detector sensitivity or from sampling/generation choices. By freezing detector thresholds, exclusion policies, sampling seeds, and generation semantics in a preregistration and archiving the validator traces, our experiment complements the literature by making the audit contract explicit and by demonstrating (artifact‑grounded) how deterministic sampling and shared generation semantics can collapse a paired confirmatory test into a degenerate outcome. This reconciliation---linking verifier traces, provider call accounting, and the detector workflow to the resulting paired contingency---is the primary methodological contribution that complements prior empirical and system‑level findings [https://openalex.org/W7157506260, https://openalex.org/W7162218766, https://openalex.org/W4404602270].

\section{Methodology}
Overview and preregistration. The experiment was preregistered (study\_id = contamination-audit-netops-2026-v1) and froze the confirmatory estimand, exclusion rules, detectors, replacement policy, sampling plan, and analysis executor prior to observing any model outputs. The primary estimand is the paired success‑rate difference (guarded - baseline) across item\ensuremath{\times}model pairs. The analysis executor is paired\_binary\_analysis\_v1 and the preregistration specifies an exact McNemar two‑sided primary test plus a paired bootstrap percentile 95\% confidence interval (10,000 resamples, bootstrap\_seed = 1729). Planned confirmatory accounting: task\_count = 150 items, planned\_episode\_count = 300 episodes (baseline + guarded), maximum\_model\_calls = 300.

Sampling and deterministic generation. To ensure reproducibility the preregistration allowed deterministic synthetic task generation. The executed run used deterministic\_netops\_task\_generator\_v5 to produce the confirmatory sample of 150 items (task\_count = 150; execution artifact task\_manifest.jsonl records per‑task generation\_seed values). Each generated task includes a source\_identifier and task\_payload that encode initial\_state, expected\_final\_state, allowed\_touched\_fields, explicit required\_sequence, and a textual intent. The preregistration required fixed seeds and deterministic selection for replacements; the execution artifacts (execution/task\_manifest.jsonl and representative task entries) show per‑task generator seeds (e.g., generation\_seed: 1..N in representative examples) and the source\_identifier prefix "synthetic-netops:" indicating provenance‑stable synthetic items in this run.

Generation semantics and effect on independent trials. The preregistration permitted both independent generation per condition and shared generation semantics; it specified independent\_condition\_generation as a parameter. The executed model\_configuration.json sets independent\_condition\_generation = false (shared\_initial\_candidate semantics). Under shared\_initial\_candidate semantics the pipeline performs a single model generation per item and reuses that shared candidate for both baseline and guarded episodes unless an exclusion/replacement rule is triggered. This execution choice reduces hosted model calls (executed model\_calls\_used = 150) relative to the preregistered maximum (300) and mechanically collapses baseline vs guarded discordance when no exclusions occur. The execution manifest and deterministic\_reconciliation artifacts record these counts (planned\_episode\_count = 300, completed\_episode\_count = 300, model\_calls\_used = 150) and confirm that provider calls were counted under a single shared generation stage per item.

Contamination detectors (implementation and thresholds). The preregistered detector suite implemented three deterministic heuristics:
1) Exact normalized token equality: canonicalize whitespace, normalize tokens, and flag exact matches between item text and archived public corpus fingerprints (high confidence on equality).
2) Fuzzy 5‑gram shingle overlap: compute 5‑gram shingles on canonicalized tokens and measure overlap/Jaccard; preregistered thresholds flagged high confidence when overlap >= 0.80, and lower confidence brackets at 0.65 and 0.50 for sensitivity analyses.
3) Provenance timestamp heuristics: compare item timestamps and DOI/URL metadata to public archive timestamps; flag items that predate the model public cutoff per heuristics in the preregistration.

Guarded policy and replacements. The preregistered guarded policy deterministically excludes only high‑confidence flagged items (exact match OR 5‑gram overlap >= 0.80). Excluded items are replaced by deterministic retemplating and topology augmentation drawn from the same task generator with fixed augmentation seeds; replacements must preserve intent similarity by an embedding safeguard (preregistered minimal cosine similarity threshold: cosine >= 0.85) before acceptance. Replacement artifacts, retemplating traces, and embedding similarity checks are included in the preregistered transformation\_manifest and were implemented in the execution pipeline; because no high‑confidence flags occurred, no replacements were performed in this run (analysis/contamination\_summary.csv empty). The preregistration additionally specified that medium‑confidence flags would receive separate handling, but the primary confirmatory exclusion mechanism applied only to high‑confidence flags.

Detector logging and artifact persistence policy. To bound artifact size the detector workflow was implemented to persist detailed per‑item fuzzy overlap scores and full overlap tables only when a threshold crossing or replacement occurred (this behavior was documented in the preregistration detection\_procedures). The rationale was to minimize unnecessary storage for large corpora while ensuring forensic traces for flagged events. As a result, in executions where no items met high‑confidence thresholds the archive contains flagless contamination\_summary artifacts but not a complete global table of per‑item overlap scores. This implementation choice is recorded in the execution transformation manifest and explains why post‑hoc aggregate overlap summaries are absent in the archive when no threshold crossings occurred.

Verifier and scoring rule (deterministic oracle). The verification oracle is deterministic\_netops\_validator\_v5 (validator\_version = 5.0). For each candidate configuration the validator performs: parsing of the configuration commands into structured updates; simulated application to an explicit device/state model (parse\ensuremath{\rightarrow}state\_update); and property checks against the task's required\_final\_state and workflow constraints (reachability/preservation rules, protected interfaces, group constraints). The preregistered scoring rule full\_intent\_constraint\_validation requires both (a) the final modeled state matches the expected\_final\_state on the allowed fields and (b) workflow constraints (sequence and transient constraints such as "must be down for field X") are respected. The validator emits binary pass/fail and detailed traces per episode: parsed\_command\_count, required\_command\_count, observed\_touched\_field\_count, violation\_codes, and normalized\_configuration. These trace fields are archived for every episode and were the basis for the binary outcomes used in analysis (execution/scoring/*.json).

Analysis plan and confirmatory executor. The preregistration specified paired\_binary\_analysis\_v1 as the primary analysis executor. For the primary estimand (paired\_success\_rate\_difference\_guarded\_minus\_baseline) we use an exact McNemar two‑sided test on the paired contingency table (n11,n10,n01,n00) and report a paired bootstrap percentile 95\% confidence interval with 10,000 resamples (bootstrap\_seed = 1729) as a descriptive companion. Secondary prespecified sensitivity analyses include: (i) treating failed model calls as failures (complete\_pair\_primary\_with\_failure\_as\_zero\_sensitivity), (ii) fuzzy‑threshold sensitivity (thresholds = 0.50, 0.65, 0.80), and (iii) subgroup analysis stratified by templated vs realistic items. The preregistration also contains power guidance showing that a two‑call per item (independent generations) or a multi‑model design increases power to detect modest paired differences (\textasciitilde{}4--6 percentage points under conservative assumptions) compared to the single‑model shared‑generation regime executed here.

Missingness, retries, and execution accounting. The preregistration specified a single automated model/API retry within 5 minutes for transient provider failures; if a retry failed the analysis plan treated the pair per the missingness contract and included a sensitivity scenario treating failed calls as failures. Execution logs show zero failed episodes and provider call accounting confirms hosted\_model\_call\_event\_count = 150, completed\_hosted\_model\_call\_event\_count = 150, failed\_hosted\_model\_call\_event\_count = 0 (deterministic\_reconciliation artifact). All seeds, provider traces, and validator traces were archived per the preregistered artifact list.

Implementation provenance and reproducibility. Key implementation artifacts frozen by preregistration and produced in execution include: execution/model\_configuration.json (declared\_model\_version = "gpt-5-mini", independent\_condition\_generation = false, max\_output\_tokens = 2000), execution/task\_manifest.jsonl (deterministic generator seeds and task\_payloads), execution/raw\_results.jsonl (shared and condition responses), and analysis artifacts (analysis/paired\_contingency\_table.csv, analysis/contamination\_summary.csv). The preregistration and execution manifest SHA‑256 values are recorded in the execution manifest and Disclosure Statement. The codepaths for detection, replacement, and validation are deterministic and archived; they can be re‑run against different sampling or generation semantics (e.g., independent\_condition\_generation = true) to increase audit sensitivity without changing detectors or statistical estimands.

\section{Results}
Execution accounting and reconciliation. The archived execution manifest records planned\_episode\_count = 300 (150 items \ensuremath{\times} 2 conditions) and completed\_episode\_count = 300 with failed\_episode\_count = 0. Because the run used shared\_initial\_candidate semantics and no high‑confidence contamination exclusions occurred, the pipeline consumed model\_calls\_used = 150 (one shared model generation per item) rather than two independent generations per condition; provider call auditing corroborates hosted\_model\_call\_event\_count = 150, cache\_miss\_count = 150 and cross\_condition\_cache\_key\_reuse\_observed = false. The deterministic\_reconciliation artifact confirms complete\_pair\_count = 150 and internal accounting consistency (n\_11 = 150, n\_10 = n\_01 = n\_00 = 0).

Contamination detectors and archived evidence of absence. Under the preregistered high‑confidence rule (exact match OR fuzzy 5‑gram overlap >= 0.80) the detector workflow produced zero high‑confidence flags for the executed sample. The analysis artifact contamination\_summary.csv is empty (no flagged rows), and the run therefore performed no deterministic exclusions or replacements. Per the preregistered detector implementation, detailed per‑item fuzzy overlap scores and replacement traces are persisted only when threshold crossings or replacements occur; because no items met the high‑confidence threshold, the archive contains the detector's flag events (none) but not a full per‑item overlap table. We therefore report the artifact‑grounded null finding (no high‑confidence flags) and emphasize that the absence of persisted overlap summaries is itself an execution artifact with diagnostic meaning (see below).

Validator outcomes: concordant passes and formal statistics. The deterministic\_netops\_validator\_v5 returned binary pass for every baseline episode (baseline\_success\_count = 150/150, baseline\_success\_rate = 1.0) and for every guarded episode (guarded\_success\_count = 150/150, guarded\_success\_rate = 1.0). The paired contingency table is n11 = 150, n10 = 0, n01 = 0, n00 = 0. The primary estimand (paired\_success\_rate\_difference\_guarded\_minus\_baseline) = 0.0 percentage points. The exact McNemar two‑sided test on these counts yields p = 1.0. The preregistered paired bootstrap percentile CI (10,000 resamples, seed = 1729) is [0.0, 0.0]. The zero‑width CI and p = 1.0 are a direct consequence of the absence of any discordant pairs; these are artifact‑verified outcomes derived from the archived paired\_contingency\_table.csv and condition\_summary.csv.

Representative diagnostic episodes. Representative archived episode traces illustrate validator behavior and the nature of shared generations. For example, pair‑000001 (task‑000001) has normalized\_response identical to the reference, parsed\_command\_count = 4 and required\_command\_count = 4, observed\_touched\_field\_count = 3, and validator verdict = VALID\_CONFIGURATION. Pair‑000002 (task‑000002) shows parsed\_command\_count = 2 (required\_command\_count = 2, level = hard) and validator verdict = VALID\_CONFIGURATION. Pair‑000003 (task‑000003) shows parsed\_command\_count = 6 (required\_command\_count = 6, level = hard) and validator verdict = VALID\_CONFIGURATION. These exemplar traces (archived per‑episode validator\_trace records) demonstrate that when the model output aligns with the preregistered intent and workflow constraints the validator produces a detailed, machine‑readable confirmation; in this execution these per‑episode confirmations are uniformly positive.

Why the paired test was uninformative for generation‑level contamination. Two artifact‑recorded execution facts explain the degenerate concordance. First, the confirmatory sample comprised provenance‑stable, deterministically generated synthetic items (source\_identifier prefix "synthetic-netops:"), so exact canonicalized matches and very high fuzzy 5‑gram overlaps to public corpus fingerprints were unlikely a priori. Second, independent\_condition\_generation = false (shared\_initial\_candidate) was in effect: when no high‑confidence flags triggered exclusions the same single model output was reused for both baseline and guarded conditions. Together these facts guarantee that, absent deterministic exclusions, baseline and guarded episodes for each item are identical in content and validator outcome---producing n11 = 150 and zero discordance by construction. The execution manifest and model\_configuration artifact explicitly encode these settings and thus ground this diagnosis in archived evidence rather than post hoc speculation.

Detector persistence policy and the missing overlap distribution. The preregistered detector persisted detailed fuzzy overlap scores only when an item crossed a threshold or when replacements were enacted; the archived absence of persisted per‑item overlap scores therefore reflects the detector's control flow under the executed conditions (no threshold crossings). This design choice reduced the post‑hoc ability to report summary overlap statistics (min/median/percentiles) for the executed sample from the archive. The missing overlap distribution is an execution artifact: it indicates that no items came close enough to the high‑confidence threshold to invoke persistence logic, but it does not reveal whether many items would have had moderate overlaps below the threshold (the archive contains no full table of low‑confidence overlaps). This factual gap motivates our concrete recommendation to persist per‑item overlap scores regardless of threshold crossings in future audits to enable transparent reporting of the overlap distribution and sensitivity analyses.

Sensitivity to preregistered design choices. The preregistration explicitly documented alternative, higher‑sensitivity designs: (a) independent\_condition\_generation = true so baseline and guarded generations are independent (doubling model calls to 300) and allowing per‑item discordance to appear even without deterministic flagging; (b) multi‑model replication to reveal model‑specific memorization signals; and (c) lower fuzzy thresholds or paraphrase‑robust detectors plus human adjudication. The archived power guidance in the preregistration explains that a 300 independent‑call design or a multi‑model plan would materially lower the detectable paired difference (planned target \ensuremath{\approx}5 percentage points at 80\% power under conservative assumptions). Because the executed run used the lower‑sensitivity shared‑generation regime and synthetic sampling, the artifact‑grounded conclusion is limited to this regime: deterministic detectors did not trigger exclusions and validators returned universal passes for the shared outputs.

Summary of archival provenance supporting the numerical results. All numeric summaries above---completed\_episode\_count = 300, model\_calls\_used = 150, baseline\_success\_count = 150, guarded\_success\_count = 150, n11 = 150, paired bootstrap CI [0.0,0.0] (10,000 resamples, seed = 1729), and McNemar p = 1.0---are computed by the preregistered analysis executor and are recorded in analysis/results.json and in the small CSV artifacts condition\_summary.csv and paired\_contingency\_table.csv. The contamination\_summary.csv artifact contains no flagged rows for this execution. Per‑episode raw responses and validator\_trace objects are archived for independent verification of every binary outcome used in the paired analysis.

\section{Discussion}
Why the result is degenerate: artifact‑grounded diagnosis. Two archived execution facts explain the degenerate null outcome. First, the sampled items were provenance‑stable synthetic tasks generated by deterministic\_netops\_task\_generator\_v5 (source\_identifier prefix "synthetic-netops:"), so canonicalized exact matches to public corpus fingerprints and high fuzzy 5‑gram overlaps above the preregistered high‑confidence threshold were unlikely. Second, execution semantics set independent\_condition\_generation = false (shared\_initial\_candidate), so when no exclusions were triggered a single shared generation per item was reused for baseline and guarded episodes. Both facts are recorded in the archived model\_configuration and task\_manifest artifacts and together produced identical baseline and guarded outputs per item, making the paired comparison uninformative for generation‑level contamination detection in this run.

What this execution answers and what it does not. The archived evidence answers the preregistered methodological question for the executed sampling regime: given provenance‑stable synthetic sampling and conservative deterministic detectors, do preregistered high‑confidence flags arise and do deterministic guarded exclusions alter validator pass rates? For this sample and settings the answer is no. The result does not generalize to web‑sourced public benchmarks, paraphrase‑sensitive contamination patterns, independent re‑generation semantics, other LLM families, or multi‑model designs. We therefore avoid extrapolating beyond the archive‑supported conclusions.

Audit sensitivity, power, and practical tradeoffs. The preregistration documented that a two‑generation per item design (independent\_condition\_generation = true) or a multi‑model plan increases the number of independent trials and therefore the power to detect paired differences of modest size (target \ensuremath{\approx}5 percentage points under conservative assumptions). The executed single‑model, shared‑generation regime reduces independent variance and makes detection of generation‑level contamination implausible unless contamination flags arise deterministically. Practically, contamination audits that aim to detect web‑sourced leakage should prioritize provenance‑rich sampling, independent generations per condition, and multi‑model replication to increase observable discordance and power.

Concrete artifact‑grounded recommendations. Based on preregistration, execution artifacts, and the deterministic reconciliation, we recommend: (1) sample provenance‑rich public items (URLs/DOIs and timestamps) to exercise detector sensitivity to real‑world leakage; (2) set independent\_condition\_generation = true so baseline and guarded conditions are produced independently and cannot collapse to shared outputs; (3) expand to multiple model instances and independent seeds to reveal cross‑model contamination signals; (4) persist per‑item fuzzy overlap scores regardless of flag status to enable overlap distribution reporting; and (5) add paraphrase‑robust detectors and human adjudication for borderline cases. Each recommendation follows from artifacts: the execution manifest and deterministic\_reconciliation confirm the current run's settings and limits, and the empty contamination\_summary motivates broader sampling and independent generation in future audits.

Operational implications for NetOps. For practitioners, the audit demonstrates that audit sensitivity depends critically on sampling provenance and generation semantics. Verifier‑backed validators are valuable oracles for functional correctness checks, but auditor designs must ensure samples exercise detector capabilities and record detector outputs at sufficient granularity for post‑hoc analysis. In production‑risk settings, independent generation, cross‑model checks, and human adjudication remain advisable.

\section{Limitations}
\begin{itemize}
\item Single‑model execution (gpt‑5-mini) --- results do not generalize across other LLM families, sizes, or training corpora.
\item Sampled items were provenance‑stable synthetic/deterministically generated tasks (source\_identifier prefix "synthetic-netops:") for this run; absence of high‑confidence flags here may not reflect contamination prevalence in real‑world, web‑crawled benchmarks.
\item Generation semantics used shared\_initial\_candidate across baseline and guarded conditions (independent\_condition\_generation = false), producing identical model outputs when no exclusions occurred and thus limiting sensitivity to interventions that rely on independent re‑generation.
\item Contamination detection relied on preregistered conservative heuristics (exact match and fuzzy 5‑gram overlap thresholds); subtler contamination patterns (paraphrased memorization, partial leakage) may escape these heuristics and require richer provenance or human adjudication.
\item Passing the deterministic\_netops\_validator\_v5 indicates satisfaction of the checked functional constraints but does not imply comprehensive production readiness or absence of semantic regressions outside the validator's scope.
\end{itemize}

\section{Conclusion}
We executed a fully preregistered, artifact‑anchored paired audit of conservative contamination detectors for LLM‑for‑NetOps evaluation. In the reproducible 150‑pair synthetic sample we observed zero preregistered high‑confidence contamination flags and identical validator pass rates (150/150 baseline, 150/150 guarded). The degenerate null outcome is artifact‑supported and stems from provenance‑stable synthetic sampling and shared generation semantics; it does not imply absence of contamination in web‑sourced benchmarks or failure of the preregistered detectors in different sampling regimes. We provide an explicit preregistration\ensuremath{\rightarrow}execution reconciliation, confirm deterministic reconciliation checks passed, and recommend concrete design changes---provenance‑rich sampling, independent generation per condition, multi‑model scope, paraphrase‑sensitive detectors, and matched power calculations---to increase audit sensitivity in future studies. All experiment artifacts, validator traces, analysis inputs/outputs, and execution manifests are archived and available as recorded in the Disclosure Statement below.

\section*{Disclosure Statement}
This study was executed under the autonomous‑run preregistration contamination-audit-netops-2026-v1 (preregistration SHA‑256 recorded in the archived bundle). LLM used: gpt-5-mini (declared\_model\_version = "gpt-5-mini"). Orchestration adapter: hosted\_netops\_gvr\_v1. Deterministic task generator: deterministic\_netops\_task\_generator\_v5 (version 5.0). Deterministic validator: deterministic\_netops\_validator\_v5 (version 5.0). Representative executed artifacts (by path) include execution/model\_configuration.json, execution/task\_manifest.jsonl, execution/raw\_results.jsonl, analysis/paired\_contingency\_table.csv, and analysis/contamination\_summary.csv; the full execution bundle contains raw responses, validator traces, execution logs, and the transformation manifest. Exact immutable initial master‑prompt reference: provenance/master\_prompt.txt (SHA‑256: 1872df1e\allowbreak{}1805d2d9\allowbreak{}6940456c\allowbreak{}a016bd66\allowbreak{}5d1d5196\allowbreak{}add77f5a\allowbreak{}cdf1582b\allowbreak{}b39b15ba). Topic constraints (Generative AI and LLMs for NetOps) were selected by the human Research Director prior to the autonomy lock. After the autonomy lock, the AI system autonomously performed literature search, gap selection, preregistration, experiment execution, analysis, manuscript drafting, peer‑review simulation, and revision. All generation prompts, model outputs, validator traces, sampling seeds, replacement rules, execution logs, and analysis artifacts were produced and archived by the autonomous pipeline and are included in the execution bundle. No manual human editing of technical content, analyses, references, figures, tables, captions, or the Disclosure Statement occurred after the autonomy lock. Pre‑lock development and rehearsal runs occurred (permitted) but pre‑lock artifacts were not used as empirical evidence for this final study unless re‑created under the locked pipeline. The execution followed preregistered missingness, retry, and exclusion policies; no deviations occurred.

\begin{thebibliography}{99}
\bibitem{ref1} https://arxiv.org/abs/2604.22513
\bibitem{ref2} https://doi.org/10.48550/arxiv.2606.06212
\bibitem{ref3} https://doi.org/10.48550/arxiv.2605.22092
\bibitem{ref4} http://citeseerx.ist.psu.edu/viewdoc/summary?doi=10.1.1.300.8659
\bibitem{ref5} https://doi.org/10.1007/s11704-026-60308-3
\bibitem{ref6} Buğra Alperen Uluırmak, Rıfat Kurban, "EvalSafetyGap: A Hybrid Survey and Conceptual Framework for LLM Evaluation-Safety Failures", 2026
\bibitem{ref7} Nguyen Van Tu, Sukhyun Nam, James Won‐Ki Hong, "Intent‐Based Network Configuration Using Large Language Models", 2024, 10.1002/nem.2313
\bibitem{ref8} Lei Huang, Weijiang Yu, Weitao Ma, Weihong Zhong, Zhangyin Feng, Haotian Wang, Qianglong Chen, Weihua Peng, Xiaocheng Feng, Bing Qin, Ting Liu, "A Survey on Hallucination in Large Language Models: Principles, Taxonomy, Challenges, and Open Questions", 2024, 10.1145/3703155
\end{thebibliography}

\end{document}
